% ============================================================
% 妙新知·逻辑客（杂志页码 34–37）
% ============================================================
\colchapter{妙新知·逻辑客}

\section*{用错地方的“放生”，与善意无关}
\auth{文/楠楠小妹}

“中央公园湖里有怪鱼出没，体型大且凶猛！”从7月中旬有消息传出“在河南平顶山市，发现一条尖嘴长牙的大鱼”开始，“湖中怪鱼”的一举一动便牵动着万千网友的心。之后长达一个多月，经过抛网捕捞、求助专业团队、抽干湖水的抓捕行动，终于有了结果：抓到了，两条，都是“幽灵火箭”鳄雀鳝。

听到这一消息，网友们激动不已。庆幸“终于拿下罪魁祸首”的有，展开联想研究“鳄雀鳝能不能吃”的有，质疑“两条鱼而已至于吗”的也有。对于种种说法和猜测，只能说一句：“怪鱼”是抓到了没错，但其背后的问题，却还没完。

鳄雀鳝，原产于北美洲。虽名为“鳝”，却绝非鳝类。鳄雀鳝有一身被菱形硬鳞覆盖的铠甲，防护力十足，且牙齿锋利，性格凶猛。它能在缺氧环境中，把空气吞到鳔中呼吸，从而在陆地上短时间生存，若是野外生长，最大甚至可长到3米左右，危害性极强。

破坏环境，袭击人类，鳄雀鳝可谓罪行累累。那么问题来了，这种原产于北美洲的怪鱼，是怎么出现在中国、频繁祸害人的？有专家推测后给出答案：放生引起的外来物种入侵。

对于大众而言，“放生”这个词并不算陌生。所谓放生，其实是指一种将被人类捕捉、有可能被宰杀的动物放归大自然的行为。而一些说法中，放生的出发点是“行善”和“积德”。可以想象，放生鳄雀鳝的人大概还觉得自己做了好事，殊不知自己的一时“好心”会惹出多大的麻烦，更不会明白“有些人放生，比杀生更残忍”。

2020年6月，浙江杭州的几位大姐为了祈福孩子高考进行放生，不巧遇上“高峰期”买不到鲤鱼，便用石斑鱼代替。然而石斑鱼属海水鱼，在淡水中根本无法存活。2016年5月，青岛山寺放生池出现“龟满为患”的现象，大量乌龟挤在池子边狭小的台阶处“晒太阳”。虽然此前寺中就有提示：放生池早已饱和，却根本拦不住前来放生的人。不仅如此，在不少人眼中，放点小鱼小虾没什么可吹嘘的，放点不一样的才算足够“虔诚”。2015年的国庆假期，北京海淀区四季曙光公园突然出现了50余只红嘴相思鸟。这种终年生活在南方，从不迁徙的鸟类为何会出现在北方？原因无他，放生。只不过，它们在寒风中不仅冻得够呛，还没有吃的……

当放生逐渐走向疯狂，没人能说清——被打开的到底是通往向善的道路，还是潘多拉的魔盒。也许很多人放生的初心是积德行善，但越来越多的事实摆在眼前：盲目放生，往往没积善，也没行善。可以预见，当善意变成跟风，便会离初衷越来越远，变成闹剧。而收起“虚伪的善良”，或许就已经是在做好事了。

\begin{quoteBox}{逻辑知点}
不负责任地“放生”，与其说是好心，不如说是借践踏他者的肩膀，把自己捧上道德的高地；是打着善良的旗号，干着让人闹心的事情。这种自以为是的感动，归根结底都是自私心态作祟，当然也得不到任何具有实际意义的积极反馈。
\end{quoteBox}

\clearpage
\section*{“热度排斥症”：让理性与专注力回归}
\auth{文/颖宝}

在瞬息万变的互联网生态下，人们习惯于火急火燎地接受资讯——上一秒因为剧情跟路人辩驳，下一秒为爱豆发新专辑欢呼。即便面对不感兴趣的话题，也要发几篇相关帖子，显示自己紧跟潮流。人们迫切将资讯全都塞进脑袋，但实则什么都没记住，如今，越来越多人开始思考是否有必要追逐每一个热点，并自我调侃患了“热度排斥症”。

热度排斥症，指对当下热门事物持观望与警惕态度，相比接受别人“种草”，更愿意用自己的节奏去接触新鲜事物。微博热议话题“热度排斥症是怎么来的”下方，不少网友表达了对某些潮流资讯的看法：“有些话题一看就是营销出来的”“知道了又对生活有什么影响呢”“越火爆越不想看”……

无尽的热点与有限的注意力，本身就是一对矛盾体。微博上，实时、文娱、视频热搜榜每天分别诞生50个话题，“同城热搜榜”每天至少27个话题。而盛产话题的平台，还包括小红书、豆瓣、B站和各大短视频软件。这意味着，要想成为“不落伍”的人，每天至少需要浏览数百种讯息。“仿佛住进互联网。”这句在贴吧论坛盛行时期就出现的梗，如今照进了现实。于是，越来越多人，开始摆脱热点的束缚。

一朝被忽悠，十年怕热搜。见太多次“网红翻车”后，大家终于意识到，有时，就算是自愿入坑，也难以从个体视角看清全貌、辨别真伪。例如，今年4月，有博主发布了一则名为“重庆也有一个北海”的短视频，上了热搜。视频中，“海水”被长条形堤坝一分为二，一半是平静的海面，一半是翻滚的浪花。站在堤坝尽头，仿佛身处大海中央。然而，质疑声很快出现——原来，此“北海”位于重庆北碚嘉陵江大桥下，是一条江，并非视觉效果上的大海。被吸引前去打卡的网友表示，堤坝颇短，并不像视频中那般“通往天边”。而在此之外，云南澄江的“粉红沙滩”、三亚清水湾的“蓝房子”等被短视频推火的景点，早前也无一例外揭下了滤镜。

影视剧领域，更是“翻车重灾区”。未播先夸、未播先评分和虚假营销，让人迷惑。2021年8月，电影《盛夏未来》被观众指责“诈骗”。预热海报中，张子枫和吴磊坐在教室里四目相对、暗藏情愫，宣传语为“希望我们可以更勇敢”，暗示着这是一部校园爱情片。显然，剧组方想吃“CP红利”，网友的讨论欲和期待值，也大多因此而起。但最终，这部“爱情片”货不对板，男女主角之间根本没有爱情戏，也导致影片口碑一度崩盘。

“让子弹飞一会儿。”好奇心再强的人，也不敢贸然试水了。对“热”趋于理性，本质上是为重新掌控情绪、喜好与审美话语权，拒绝跟风，找回真正的精神支柱。或许，只有等没人再讨论“热度排斥症”时，我们才算真正地回归理性。

\begin{quoteBox}{逻辑知点}
“按头种草”“热搜榜”这样的推送出现之后，人们很容易生出一种被强迫的感觉，毕竟“落伍”的隐忧真实而具体，大家不得不去追寻热度，以维持自己在现代社会中的节拍。然而，漫无目的地“追热点”“刷热搜”恰恰造就了热点潮流对人的“绑架”，我们的兴趣爱好、社交方式、消费习惯等等都应留待自己去寻找并抉择。
\end{quoteBox}

\clearpage
\section*{没有沟通的“天台”，少年该向谁说？}
\auth{文/柳飘飘}

最近，几张令人窒息的截图，又在各个社交媒体流传：一个小孩站在天台上，控诉自己的父亲，因为自己一笔一笔写出来的56本小说，被父亲一页一页地撕掉了……

这些截图，来自2018年开播的《少年说》。你没看过？不奇怪。因为它一开播便恶评如潮，甚至涉嫌抄袭。而奇怪的是，四年过去了，它非但没凉，热度还越来越高。或许，最真实的中国式亲子关系，最真实的青少年群体，不在那些剧里，就在这部口碑如此分裂的综艺里。《少年说》的策划出发点，无疑是好的。这一代少年极度渴望直接地表达和沟通，然而父辈都是不善于面对面沟通、直白表达情绪的人。于是，节目组让孩子们站上学校天台，向同学朋友，向自己的爸妈说出平时难以启齿的，或没机会说的真心话。

只不过，这个综艺从2018年播到2020年，两代人的距离，真的拉近了吗？节目中，孩子们多是抱着能打破和父母之间的壁垒，进行一次有效沟通的初衷来的。而底下的父母面对孩子真诚的心声，却往往是继续打压。就像文章开篇说到的小女孩，她从小就喜欢写小说，文科成绩也很好。小女孩这样形容自己50多本小说被撕毁后的感觉：“那段时间，我就像失去了灵魂一样。”可台下的爸爸听到小孩这样剖析自己受伤的心灵后，继续“插刀”：“先把你的小说放一放，等我们的数学成绩提上去，我们再商定一下。”女孩勇敢地走上天台，是希望爸爸能尊重她，可爸爸完全没理解她的意思。

这样受伤的孩子还有很多。在这个节目里，你可以观察到中国式父母一个奇特的现象：当父母终于听到孩子的心声时，他们的第一反应往往是为自己“找台阶下”，而不是去体会孩子为什么这样说。父母和孩子身处的物理高度虽然互换了，可一旦沟通起来，站在天台上喊话的孩子，不过是换了一个位置，在听父母的教训。

为什么四年间这个节目仍在不断出圈？其实功劳就在这一个个真实的孩子身上。节目中，每一个孩子的表达都很真诚，他们真的把它当成一次和父母沟通的好机会，才选择勇敢地走上天台，掏心掏肺，把自己的委屈通通倒出来。不过，“天台喊话”并没能解决两代人的隔阂，反而暴露了这种隔阂的无法弥合。“少年说”，变成了“说了也白说”。

那么，四年之后，我们的少年有更好的表达自我的渠道了吗？依然没有。不管从节目、视频还是生活中，你都很难看到像《少年说》里的少年们那样真实的内心抒发。他们的内心越封闭，形象就越模糊，渐渐地，他们不再说话，也没人去真正关心他们在想什么。基于此，也就不难理解《少年说》时不时的爆火了——它虽然不够好，但再也没有一款和它相似的，甚至比它更好的节目。

\begin{quoteBox}{逻辑知点}
纵观《少年说》里的孩子，都非常懂事。与其说他们是在“控诉”，不如说是怀着忐忑的心情去跟父母提一点意见；长辈回话的时候，他们也认真倾听，反思自己……而纵然如此沉重，孩子们的心声仍然淹没在千千万万父辈权威所筑起的高墙里，需要更多的平台与机会去为他们发声、共情。
\end{quoteBox}

\clearpage
\section*{“被胁迫直播”：闹剧散场，透支同情}
\auth{文/王动}

在电视剧的片头，往往都有一段声明：“本故事纯属虚构，如有雷同，实属巧合。”但与不断强调自己虚构性的影视剧不同，直播处处暗示自己的真实性，观众也愿意相信，他们隔着屏幕看到的，是世界另一个角落里真切发生着的事情。不辜负这样一份信赖，努力维持诚信原则，就是主播与观众的默契，也是直播行业持久发展的不二选择。

人人都听过《狼来了》的故事。可没想到，2022年的互联网上，寓言照进了现实。“如果需要报警，你就摸一摸鼻子。”8月6日，主播“大嘴燕燕”的直播间，有人对主播留言说。原以为是网友多虑了，可主播随后就对着镜头，摸了摸鼻子。屏幕前的网友顿感毛骨悚然，而越来越多的人聚集在了这个每天准时开播的直播间。人们逐渐发觉，这个直播间，简直处处透着诡异。

最初令人感到不对劲的，是主播的直播环境。一个“正常”的直播间通常是经过精心布置的，背景布、补光灯都是常备配件。但是“大嘴燕燕”直播的地方，看起来却十分简陋——那里似乎是一座毛坯平房，墙面上凌乱地分布着裸露的电线，窗外则罩上了防盗网……尽管出镜的只有主播，但房间内的种种痕迹，暗示了另一个人的存在：窗户上挂着一件衬衫，显然是男式的，远处的桌子上有一张结婚照，但上面的新娘和主播看起来似乎并不是同一个人。主播更像是不小心闯入这间屋子的外来者，或者，换句话说，主播像是被强行关在了这个不属于她的房间里。

当然，“诡异”的还有主播本人。虽然总是面带微笑，但她的动作、表情却很僵硬，也从不跟留言区的观众有任何语言交流。就此，直播间的网友嗅到了危险的气息，有人试着用某些“暗号”和主播交流：“如果被绑架了，就眨一眨眼睛”“如果需要报警，就摸一下耳垂”……向来跟观众没什么互动的主播，这时却对网友们抛来的暗号反应迅速，一一照做。

难道说，主播真的被胁迫了？不少人都对此深信不疑。抱着“不要漏掉一个受害者”的态度，有热心网友还是报了警。最终，在“大嘴燕燕”事件登上热搜之后的第二天，警方就给出了调查结果。之所以能这么快，是因为这起事件其实非常简单：“该女主播并未受到任何胁迫，是为了增加个人账号粉丝，引起网友关注，进行了恶意误导。”并没有人在这起闹剧中受到真实的人身伤害，但坏消息是，“狼来了”的名额，又被白白消耗了一次。

其实，假装被胁迫这种套路，早在直播行业兴起初期，就已经出现。而为了火，直播里上演的大戏恐怕远不止“被胁迫”这一幕。通过编织虚假的苦难与困境，某些“网红”骗取了同情心与关注度，也就抓住了所谓的“流量”。可这些人往往忽视了最重要的一点：“同情心”是一种不可再生资源。通过欺骗得来的流量，最终会全部还回去，同情心和信任感的崩塌，同样难以补救。

在“大嘴燕燕”的新闻下，有人评论道：宁愿被耍一百次，也不想错过一次帮助受害人脱离苦海的机会。这当然很了不起。但每一次同情心的虚耗，最终都会削减这个社会的善意。“狼来了”的把戏玩得太多，哪天狼真的来了，我们每个人都可能成为受害者。

\begin{quoteBox}{逻辑知点}
当人们把“同情心”当作一种可以反复消费的流量密码，每一次虚假的苦难叙事都会让真正的求助变得可疑。被透支的信任无法轻易修复，而一个社会的善意，恰恰建立在一次次真实的被回应之上。守住诚信的底线，也就是守住人间最基本的温度。
\end{quoteBox}
